% Clase 3 --- Dipolo sobre el plano bisector: los campos de +Q y -Q en P tienen
% igual módulo; las componentes horizontales se cancelan y el neto apunta en -Y.
\begin{center}
\begin{tikzpicture}
  % ejes
  \draw[figaux] (0,-1.9) -- (0,1.9);
  \draw[figaux] (-0.4,0) -- (3.4,0);
  % cargas
  \node[figpos] at (0,1.3) {};
  \node[figlbl,left=4pt] at (0,1.3) {$+Q$};
  \node[figneg] at (0,-1.3) {};
  \node[figlbl,left=4pt] at (0,-1.3) {$-Q$};
  % separación d
  \draw[figaux] (-1.6,1.3) -- (-1.6,-1.3);
  \node[figlblaux,left=2pt] at (-1.6,0){$d$};
  % punto P
  \node[figneutra, minimum size=4pt] at (2.6,0) {};
  \node[figlbl,above right=1pt] at (2.6,0){$P$};
  \node[figlblaux,above=2pt] at (1.3,0.02){$x$};
  % líneas r
  \draw[figgray,line width=0.6pt] (0,1.3) -- (2.6,0);
  \draw[figgray,line width=0.6pt] (0,-1.3) -- (2.6,0);
  % campos en P
  \draw[figvec,figred] (2.6,0) -- (3.59,-0.49);
  \node[figlbl,figred,right=1pt] at (3.59,-0.49) {$\vec E_{+}$};
  \draw[figvec] (2.6,0) -- (1.61,-0.49);
  \node[figlbl,figblue,left=1pt] at (1.61,-0.49) {$\vec E_{-}$};
  \draw[figvecres] (2.6,0) -- (2.6,-1.05);
  \node[figlbl,figred,right=2pt] at (2.6,-0.75) {$\vec E$};
\end{tikzpicture}\\[0.4em]
{\small\itshape En el plano bisector, las componentes horizontales de
$\vec E_{+}$ y $\vec E_{-}$ se cancelan; el campo neto $\vec E$ apunta en
$-\hat{Y}$ (de la carga positiva a la negativa).}
\end{center}
